\documentclass[a4paper,11pt]{article}
\usepackage[utf8]{inputenc}
\usepackage[catalan]{babel}
\usepackage{amsmath,amssymb}
\usepackage{booktabs}
\usepackage{geometry}
\geometry{top=2.2cm,bottom=2.2cm,left=2.5cm,right=2.5cm}
\usepackage{microtype}
\usepackage{parskip}
\setlength{\parskip}{5pt}
\setlength{\parindent}{0pt}

\title{\textbf{Control híbrid MPC\,+\,RL per a sistemes d'emmagatzematge\\
en comunitats energètiques}\\[5pt]
\large Resum tècnic}
\author{David Bruguera Tornés \quad|\quad Tutor: Andrés El-Fakdi Sencianes\\
\small Màster en Ciència de Dades --- Juny 2026}
\date{}

\begin{document}
\maketitle
\thispagestyle{empty}

% ============================================================
\paragraph{Cas d'estudi.}
Comunitat energètica a Cornellà de Llobregat: 20~habitatges i una escola
amb una instal·lació fotovoltaica compartida de 15~kW de pic i un BESS
(\emph{Battery Energy Storage System}) de 50~kWh / 50~kW d'inversor.
El simulador de la comunitat genera 7.300~dies de dades de generació solar,
demanda agregada i preus PVPC;
es divideixen en 80\,\% d'entrenament (5.840 dies) i 20\,\% de test (1.460~dies).

% ============================================================
\paragraph{Notació.}
\begin{center}\small
\begin{tabular}{@{}ll@{}}
\toprule
\textbf{Símbol} & \textbf{Significat} \\
\midrule
$t,\,\tau$                        & instants de temps; $\Delta t = 1\,\text{min} = 1/60\,\text{h}$ \\
$P_t^{\mathrm{batt}}$             & potència bateria (kW); positiu = càrrega, negatiu = descàrrega \\
$P_{\max}$                        & potència màxima de l'inversor (50\,kW) \\
$\mathrm{SoC}_t$                  & estat de càrrega [0,\,1]; capacitat $C_{\mathrm{bat}}=50$\,kWh \\
$\eta$                            & eficiència de la bateria ($\eta=1$ en la implementació) \\
$E^+_t,\;E^-_t$                   & energia comprada / venuda a la xarxa (kWh, $\geq 0$) \\
$p^+_t,\;p^-_t$                   & preu de compra / venda de l'electricitat (€/kWh) \\
$D_t,\;G_t$                       & demanda total i generació solar (kW) \\
$c_{\mathrm{deg}}$                & cost de degradació de la bateria (0{,}010\,€/kWh) \\
$\pi_\theta,\;s_t,\;a_t$          & \textit{policy} RL, vector d'estat ($\mathbb{R}^8$) i acció \\
$r_t,\;\gamma,\;\hat{A}_t$        & recompensa, factor de descompte i avantatge estimat (GAE) \\
$\mathrm{SoC}_{\mathrm{ref}}$     & SoC objectiu proposat per l'agent RL (MPC+RL) \\
$w=20$                            & pes de la penalització de seguiment del SoC (MPCTracker) \\
\bottomrule
\end{tabular}
\end{center}

% ============================================================
\section*{1.\ Model de la bateria}

L'estat de càrrega (\emph{State of Charge}, SoC) evoluciona a cada pas
de $\Delta t = 1/60$\,h (1~minut) segons:
\begin{equation}
  \mathrm{SoC}_{t+1} = \mathrm{SoC}_t
  + \frac{\eta \cdot P_t^{\mathrm{batt}} \cdot \Delta t}{C_{\mathrm{bat}}},
  \qquad \eta = 1
  \label{eq:batt}
\end{equation}
amb $C_{\mathrm{bat}}=50$\,kWh i les restriccions físiques
$P_t^{\mathrm{batt}}\in[-50,\,50]$\,kW i $\mathrm{SoC}_t\in[0,1]$.
L'eficiència unitària ($\eta=1$) simplifica la dinàmica; addicionalment,
el cost de degradació es dobla si el SoC opera fora del rang
$[0{,}10,\;0{,}90]$ per penalitzar l'ús als extrems.

% ============================================================
\section*{2.\ Controladors}

\subsection*{MPC($H=1$) — greedy d'un sol pas}

En cada minut $t$, resol el problema d'optimització convexa ($H=1$):
\begin{equation}
  \max_{\{P_\tau,\,E^+_\tau,\,E^-_\tau\}}
  \sum_{\tau=t}^{t+H-1}
  \Bigl[(p^-_\tau \cdot E^-_\tau) - (p^+_\tau \cdot E^+_\tau)
        - (c_{\mathrm{deg}} \cdot |P_\tau| \cdot \Delta t)\Bigr]
\end{equation}
\begin{align*}
  \text{subjecte a}\quad
  & E^+_\tau - E^-_\tau = (D_\tau - G_\tau + P_\tau) \cdot \Delta t
    \quad (E^+_\tau,\,E^-_\tau \geq 0)\\
  & -P_{\max} \leq P_\tau \leq P_{\max},\quad
    \mathrm{SoC}_{\min} \leq \mathrm{SoC}_{\tau+1} \leq \mathrm{SoC}_{\max}
\end{align*}
S'optimitzen conjuntament la potència de la bateria ($P_\tau$) i els fluxos
d'energia amb la xarxa ($E^+_\tau$, $E^-_\tau$): es maximitza el que es ven
($p^-\cdot E^-$, ingressos) menys el que es compra ($p^+\cdot E^+$, cost) i
la penalització de degradació per ciclatge.
La restricció de balanç assegura que el flux net amb la xarxa cobreix
exactament la demanda neta de la comunitat ($D_\tau - G_\tau$) més la
potència de la bateria; les restriccions de SoC i potència garanteixen
l'operació física segura.
Implementat amb CVXPY i el \textit{solver} CLARABEL.
Sense horitzó, no pot anticipar períodes favorables: com que el preu de
venda (0,080\,€/kWh) és inferior al de compra (0,120\,€/kWh), el MPC($H=1$)
opta sistemàticament per la \textbf{inacció} (100\,\% del temps),
constituint la cota inferior del sistema.

\subsection*{MPC($H=60$) — horitzó de 60~minuts}

Estén el mateix problema a $H=60$ passos amb coneixement perfecte dels
preus, generació i demanda futurs, resolent un sol
LP (\textit{Linear Program}) per iteració:
\begin{equation}
  \max_{\{P_\tau,\,E^+_\tau,\,E^-_\tau\}}
  \sum_{\tau=t}^{t+H-1}
  \Bigl[(p^-_\tau \cdot E^-_\tau) - (p^+_\tau \cdot E^+_\tau)
        - (c_{\mathrm{deg}} \cdot |P_\tau| \cdot \Delta t)\Bigr]
\end{equation}
\begin{align*}
  \text{subjecte a}\quad
  & E^+_\tau - E^-_\tau = (D_\tau - G_\tau + P_\tau) \cdot \Delta t,
    \quad E^+_\tau,\,E^-_\tau \geq 0\\
  & -P_{\max} \leq P_\tau \leq P_{\max},\quad
    \mathrm{SoC}_{\min} \leq \mathrm{SoC}_\tau \leq \mathrm{SoC}_{\max}
\end{align*}
L'estructura és idèntica al MPC($H=1$) però l'horitzó de 60 passos permet
anticipar variacions de preus o generació i planificar càrregues/descàrregues
en moments econòmicament favorables.
Actua el 7,6\,\% del temps i millora el MPC($H=1$) en
324\,€ en 1.460 dies (0,22\,€/dia), benefici marginal molt petit.

\subsection*{Agent RL — PPO (\textit{Proximal Policy Optimization})}

\paragraph{\textit{Policy function} i xarxa neuronal.}
L'agent modela una \textit{policy} $\pi_\theta(a_t\mid s_t)$ parametritzada
per una xarxa neuronal MLP (\textit{Multi-Layer Perceptron}) amb funció
d'activació ReLU (\textit{Rectified Linear Unit}).
La transició entre estats segons aquesta \textit{policy} es defineix com:
\begin{equation}
  s_t\in\mathbb{R}^8
  \;\xrightarrow[{\;a_t\in[-1,1]\;}]{\;}\;
  s_{t+1}\in\mathbb{R}^8
\end{equation}
El vector d'estat $s_t$ conté 8 variables normalitzades: SoC, generació solar,
demanda escola, demanda habitatges, preu de l'electricitat i codificació cíclica
del temps ($\sin\theta_t$, $\cos\theta_t$, hora normalitzada).
L'acció $a_t=1$ correspon a càrrega màxima i $a_t=-1$ a descàrrega màxima.

L'arquitectura \textbf{actor-crític} (\textit{actor-critic}) separa la
representació en dues xarxes independents: l'\textbf{actor} aprèn la
\textit{policy} $\pi_\theta(a_t|s_t)$ (quina acció executar donat l'estat),
i el \textbf{crític} estima el valor esperat de cada estat $V(s_t)$ per
calcular l'avantatge i guiar el gradient de l'actor sense modificar-ne
la \textit{policy} directament. Actor: \textbf{1 capa de 64 neurones};
crític: \textbf{2 capes de 64 neurones}.

\paragraph{\textit{Reward function}.}
La recompensa principal és el \textbf{guany marginal respecte a no usar
la bateria}, menys la penalització de degradació:
\begin{equation}
  r_t = \underbrace{\bigl(b_t^{\mathrm{batt}} - b_t^{\mathrm{ref}}\bigr)}_{\text{guany marginal}}
       - \underbrace{\alpha_{\mathrm{deg}}\cdot c_{\mathrm{deg}}\cdot
         |P_t^{\mathrm{batt}}|\cdot\Delta t}_{\text{cost degradació}},
  \qquad c_{\mathrm{deg}}=0{,}010\,\text{€/kWh}
  \label{eq:reward}
\end{equation}
on $b_t^{\mathrm{batt}}$ és el benefici econòmic net amb bateria,
$b_t^{\mathrm{ref}}$ el benefici sense bateria i $\alpha_{\mathrm{deg}}$
és un factor calibrat experimentalment per evitar la \textit{policy} d'inacció.

\paragraph{Retorn descomptat i objectiu de PPO.}
L'agent maximitza el retorn esperat descomptat:
\begin{equation}
  J(\theta) = \mathbb{E}_{\pi_\theta}\!\left[\sum_{t=0}^{T}\gamma^t r_{t+1}\right],
  \qquad \gamma = 0{,}999
\end{equation}
PPO estima el gradient via l'avantatge generalitzat
(GAE, \textit{Generalized Advantage Estimation}, $\lambda=0{,}98$) i aplica
una restricció de \textit{clip} per garantir actualitzacions estables:
\begin{equation}
  L^{\mathrm{CLIP}}(\theta)=\mathbb{E}_t\!\left[
    \min\!\left(r_t(\theta)\,\hat{A}_t,\;
    \mathrm{clip}(r_t(\theta),1\pm\epsilon)\,\hat{A}_t\right)\right],
  \quad\epsilon=0{,}2
\end{equation}
on $r_t(\theta)=\pi_\theta(a_t|s_t)/\pi_{\theta_{\rm old}}(a_t|s_t)$
és el ràtio de probabilitats i $\hat{A}_t$ l'avantatge estimat.
Entrenament: \textbf{2.000.000 timesteps} sobre 5.840 dies d'entrenament.

% ============================================================
\section*{3.\ Arquitectura híbrida MPC\,+\,RL}

La idea central és combinar l'adaptabilitat de l'agent RL amb la capacitat
d'optimització multi-pas del MPC. L'arquitectura té dos nivells temporals:
un \textbf{macro-pas} (decisió de l'RL, cada 60 minuts) i un
\textbf{micro-pas} (execució del MPC, cada minut).

\paragraph{Macro-pas — decisió de l'RL (cada 60~minuts).}
L'agent RL observa l'estat $s_t$ i proposa un \textbf{SoC objectiu}:
\begin{equation}
  \mathrm{SoC}_{\mathrm{ref}} = \pi_\theta(s_t) \in [0,1]
\end{equation}
que la bateria hauria d'assolir en els pròxims 60~minuts.
Per a aquest agent, l'arquitectura és \textbf{[64,64]} per a l'actor i
el crític ($\gamma=0{,}994$, $\lambda=0{,}95$), i la recompensa és el
guany marginal respecte a l'escenari sense bateria
(eq.~\eqref{eq:reward}).

\paragraph{Micro-pas — execució del MPC (cada minut).}
Un \textbf{MPCTracker} amb horitzó $H=60$ (CVXPY/CLARABEL) resol en cada
minut un problema de programació quadràtica que minimitza el cost econòmic
de transaccions més una penalització per desviació respecte al SoC objectiu:
\begin{equation}
  \min_{\{P_\tau,\,E^+_\tau,\,E^-_\tau\}}
  \sum_{\tau=t}^{t+H-1}
  \Bigl[(p^+_\tau \cdot E^+_\tau) - (p^-_\tau \cdot E^-_\tau)\Bigr]
  + w\cdot\bigl(\mathrm{SoC}_{t+H_{\mathrm{rem}}} - \mathrm{SoC}_{\mathrm{ref}}\bigr)^2
\end{equation}
\begin{align*}
  \text{subjecte a}\quad
  & E^+_\tau - E^-_\tau = (D_\tau - G_\tau + P_\tau)\cdot\Delta t,
    \quad E^+_\tau,\,E^-_\tau \geq 0\\
  & -P_{\max} \leq P_\tau \leq P_{\max},\quad
    \mathrm{SoC}_{\min} \leq \mathrm{SoC}_\tau \leq \mathrm{SoC}_{\max}
\end{align*}
Es minimitza el cost net de transaccions amb la xarxa (comprar menys vendre),
penalitzant la desviació del SoC respecte al SoC objectiu proposat per l'agent.
La penalització $w\,(\cdot)^2$ s'aplica únicament al pas final del macro-pas
($t+H_{\mathrm{rem}}$), no a tots els passos de l'horitzó.
L'acció executada és $P_t^* = P_\tau[0]$ (primer pas de la seqüència òptima),
i un filtre de seguretat garanteix els límits de SoC i potència.

L'arquitectura garanteix \textbf{zero violacions de restriccions} sense
necessitat de reentrenar l'agent, i produeix els millors resultats econòmics
de tots els controladors avaluats.

% ============================================================
\section*{4.\ Resultats (1.460 dies de test)}

\begin{table}[h]
\centering
\small
\begin{tabular}{@{}lrrrr@{}}
\toprule
\textbf{Controlador} & \textbf{Benefici net (€)} &
  \textbf{Estalvi vs MPC($H$=1)} & \textbf{Cicles bat.} &
  \textbf{SoC mitjà} \\
\midrule
MPC($H=1$)      & $-97.826$ & ---                         & 0    & 0{,}000 \\
MPC($H=60$)     & $-97.502$ & $+324$\,€\; $(+0{,}3\,\%)$ & 247  & 0{,}010 \\
RL (PPO)        & $-94.035$ & $+3.792$\,€\; $(+3{,}9\,\%)$& 892  & 0{,}230 \\
\textbf{MPC+RL} & $\mathbf{-87.142}$ &
  $\mathbf{+10.684}$\,\textbf{€\; $(+10{,}9\,\%)$} &
  1.951 & \textbf{0{,}705} \\
\bottomrule
\end{tabular}
\medskip\\
\footnotesize
\textit{Avaluació sobre 2.102.400 minuts de simulació.
Temps de decisió: MPC($H=1$) 0,63\,ms; MPC($H=60$) 3,77\,ms;
RL 0,13\,ms; MPC+RL 2,06\,ms — tots compatibles amb temps real.}
\end{table}

\textbf{Conclusions principals:}

\begin{itemize}
  \item \textbf{MPC+RL} és el millor en totes les mètriques econòmiques:
    estalvi de \textbf{7,32\,€/dia} respecte a la inacció total del
    MPC($H=1$), i \textbf{4,72\,€/dia} addicionals sobre el RL en solitari.

  \item \textbf{L'agent RL} supera el MPC($H=60$) en 3.468\,€ sense tenir
    accés explícit a previsions de generació ni demanda; la seva \textit{policy}
    de ciclatge és quasi-perfectament simètrica
    (44.583 vs 44.581\,kWh en 1.460~dies).

  \item La \textbf{capa de projecció} garanteix simetria exacta de
    càrrega/descàrrega (97.527\,kWh en cada sentit) i zero violacions de
    restriccions; la seguretat operativa no requereix reentrenar l'agent.

  \item L'\textbf{estudi d'ablació} (4 activacions $\times$ 2 valors de
    GAE-$\lambda$, cadascuna 500k timesteps) confirma que el sistema és
    robust als hiperparàmetres: rang de variació $<7\,\%$ entre totes
    les configuracions, gràcies a la simplificació que la capa de
    projecció imposa sobre la tasca de l'agent.
\end{itemize}

\end{document}
